%%=============================================================================
%% DADOS DO TRABALHO -- FONTE ÚNICA DA VERDADE
%%
%% Este é o ÚNICO arquivo onde os dados do seu trabalho devem ser digitados.
%% Tudo o que aparece na capa, folha de rosto, ficha catalográfica, folha de
%% aprovação e nos metadados do PDF é gerado a partir daqui.
%%
%% Campos marcados como OBRIGATÓRIO que ficarem em branco aparecem no PDF
%% como "[[ PREENCHER: campo ]]" e geram aviso na compilação. Ao gerar a
%% versão final, use a opção `entrega' em main.tex: ela transforma o aviso
%% em erro e impede a compilação com campos pendentes.
%%
%% Para reaproveitar estes dados no corpo do texto, use os acessores:
%%   \TituloTrabalho  \NomeAutor      \NomeOrientador   \NomeCoorientador
%%   \NomeCurso       \NomePrograma   \NomeInstituicao  \AreaConcentracao
%%   \LinhaPesquisa   \LocalTrabalho  \CidadeTrabalho   \AnoTrabalho
%% Exemplo: "Ao meu orientador, \NomeOrientador, agradeço..."
%%=============================================================================

%% --- Identificação ---------------------------------------------------------
\titulo{Avaliação Comparativa de Arquiteturas RAG para Respostas Jurídico-Administrativas Confiáveis}  % OBRIGATÓRIO
\subtitulo{recuperação densa, esparsa, híbrida (denso-esparso) e por grafo (GraphRAG) com geração single-pass e multiagente no corpus jurídico-administrativo brasileiro}
\titulotraduzido{Comparative Evaluation of RAG Architectures for Reliable Legal-Administrative Answers: Dense, Sparse, Hybrid (Dense-Sparse) and Graph-Based (GraphRAG) Retrieval with Single-Pass and Multi-Agent Generation over a Brazilian Legal-Administrative Corpus}  % para o abstract

\autor{Pedro Ivo Guimarães Póvoa}                 % OBRIGATÓRIO
\autorficha{PÓVOA, Pedro Ivo Guimarães}           % OBRIGATÓRIO (formato invertido)

\orientador{Prof. Dr. Alex Lopes Pereira}           % OBRIGATÓRIO
\coorientador{}                                     % opcional

%% --- Vínculo institucional -------------------------------------------------
\curso{}                                 % não se aplica (mestrado)
\programa{Administração Pública}         % OBRIGATÓRIO em mestrado/doutorado
\area{}                                  % área de concentração (pós-graduação)
\linhapesquisa{}                         % linha de pesquisa (pós-graduação)

\cidade{Brasília}                        % usada na folha de aprovação
\local{Brasília-DF}                      % usado na capa e folha de rosto
\data{2026}                              % ano de depósito

%% --- Defesa ----------------------------------------------------------------
%% Informe APENAS a data; a cidade vem de \cidade acima.
\datadefesa{3 de maio de 2027}            % OBRIGATÓRIO -> "Brasília, 3 de ... ." (provisória; banca ainda a definir)

%% Banca: \membrobanca[Instituição]{Nome}{Função}   (instituição é opcional)
%% Banca ainda não definida formalmente pelo programa -- só o orientador é certo.
\membrobanca{Prof. Dr. Alex Lopes Pereira}{Orientador}
\membrobanca{A definir}{Examinador}
\membrobanca{A definir}{Examinador}

%% --- Ficha catalográfica ---------------------------------------------------
%% A ficha OFICIAL é elaborada pela Biblioteca Ministro Moreira Alves.
%% Solicite por biblioteca@idp.edu.br e transcreva os dados recebidos aqui.
\cutter{}                                % notação Cutter fornecida pela Biblioteca (a definir)
\numeropaginas{00}                       % OBRIGATÓRIO: total de folhas (a definir)
\cdd{XXX}                                % OBRIGATÓRIO: Classificação Decimal de Dewey (a definir)
\assuntos{1. Inteligência artificial. 2. Recuperação da informação. 3. Administração pública -- Processamento de dados.}

%% --- Resumo ----------------------------------------------------------------
\palavraschave{Recuperação Aumentada por Geração (RAG). GraphRAG. Grafos de Conhecimento. Geração Verificada. Corpus Jurídico-Administrativo.}   % OBRIGATÓRIO
\keywords{Retrieval-Augmented Generation (RAG). GraphRAG. Knowledge Graphs. Verified Generation. Legal-Administrative Corpus.}             % OBRIGATÓRIO

%% --- Sobreposições opcionais -----------------------------------------------
%% Normalmente você NÃO precisa mexer daqui para baixo.

%% Texto da natureza do trabalho. Por padrão é gerado a partir da opção da
%% classe (bacharelado/especializacao/mestrado/doutorado).
% \preambulo{Monografia apresentada como requisito para obtenção do título de
%            bacharel em Direito do Instituto Brasileiro de Ensino,
%            Desenvolvimento e Pesquisa -- IDP.}

%% Tipo do trabalho na ficha catalográfica (padrão: Monografia/Dissertação/Tese)
% \tipotrabalho{Trabalho de Conclusão de Curso}

%% Linha inteira da data de defesa, se o seu programa exigir outra redação
% \datadefesacompleta{Aprovada em 15 de dezembro de 2026.}

%% Logotipo da capa
% \logo{figuras/idp-logo.png}
% \semlogo
